\section{Case study: the soft-knee limiter}\label{sec:case}

This section proves the property that the verifier must establish for the case study, and
characterizes exactly where the first draft fails. The verifier does not exist yet, so the proofs
are by hand; an exhaustive enumeration then checks both statements mechanically.

\subsection{Floating-point facts}

Let $\Fl$ be the binary32 values of \cref{def:onenan} and
$\Om=(2-2^{-23})\cdot2^{127}$ the largest finite one. Let $\rn:\R\to\Fl\cup\{\pm\infty\}$ round to
nearest, ties to even. We use the following facts~\cite{ieee754-2019,muller2018handbook}.
\begin{enumerate}[label=(F\arabic*),ref=F\arabic*]
\item\label{f:mono} \emph{Monotonicity}: $x\le y$ implies $\rn(x)\le\rn(y)$.
\item\label{f:exact} \emph{Exactness}: if $x\in\Fl$ then $\rn(x)=x$.
\item\label{f:sterbenz} \emph{Sterbenz's lemma}~\cite{sterbenz1974}: if $a,b\in\Fl$ are finite and
  $b/2\le a\le 2b$, then $a-b\in\Fl$.
\item\label{f:ops} \emph{Operations}: for non-NaN operands, \op{FADD}, \op{FSUB}, \op{FMUL},
  \op{FDIV} and \op{FMA} return $\rn$ of the exact result (\op{FMA} with a single rounding),
  except for the invalid cases $\infty-\infty$, $0\cdot\infty$, $0/0$ and $\infty/\infty$,
  which return $\NaN$. An exact zero result of a subtraction of equal values is $+0$; $1+(+\infty)
  =+\infty$; $1/(+\infty)=+0$; a finite non-negative value divided by a positive finite value
  is never $\NaN$.
\item\label{f:overflow} \emph{Overflow}: $\rn(x)=+\infty$ exactly when $x\ge2^{128}-2^{103}$.
\end{enumerate}
Write $\kappa$ for the binary32 value of the literal \kw{0.99}:
\[
\kappa=\frac{4152361}{4194304}=0.9900000095367431640625\qquad(\texttt{0x3F7D70A4}).
\]

\subsection{The repaired kernel is correct}

\begin{theorem}[Soft knee]\label{thm:softknee}
For every $x\in\Fl$ with $x$ finite and every $\mathit{knee}\in\Fl$ with
$\frac12\le\mathit{knee}\le\kappa$, the kernel of \cref{lst:softknee} does not trap, none of its
ports holds $\NaN$, and its result $r$ satisfies $-1\le r\le1$ and
$(x\ge0)\Leftrightarrow(r\ge0)$, with IEEE comparisons.
\end{theorem}
\begin{proof}
We follow the ports in order and write $k$ for $\mathit{knee}$.
\begin{itemize}
\item $\%0=|x|\in[0,\Om]$, and $\%0=+0$ for $x=\pm0$.
\item $\%1=\rn(\%0-k)$ is finite. If $\%0>k$, the exact difference is positive and at least
  $2^{-24}$, the spacing of $\Fl$ near $k$, so $\%1>0$ by \ref{f:mono} and \ref{f:exact}; if
  $\%0=k$, $\%1=+0$ by \ref{f:ops}; if $\%0<k$, $\%1<0$.
\item $\%2=\op{FMAX}(\%1,+0)$ is $\%1$ when $\%1>0$ and $+0$ otherwise (the maximum places
  $+0$ above $-0$). So $\%2\in\{+0\}\cup(0,\Om]$; it is never $-0$ or $\NaN$.
\item $\%3=1-k$ exactly, by \ref{f:sterbenz} with $a=1$, $b=k$: $k/2\le\frac12<1\le2k$. Hence
  $\%3\in[1-\kappa,\frac12]$, and $1-\kappa=41943/4194304>0$.
\item $\%4=\%2/\%3$: a value in $[0,\Om]$ divided by a positive finite value. By \ref{f:ops}
  and \ref{f:mono}, $\%4\in[+0,+\infty]$; it is $+\infty$ when the quotient overflows and never
  $\NaN$.
\item $\%5=1+\%4\in[1,+\infty]$: if $\%4$ is finite, $1+\%4\ge1$ and \ref{f:mono},
  \ref{f:exact} give $\%5\ge1$; if $\%4=+\infty$, $\%5=+\infty$.
\item $\%6=1/\%5\in[0,1]$: if $\%5=+\infty$, $\%6=+0$; otherwise $1/\%5\in(0,1]$ and
  \ref{f:mono}, \ref{f:exact} give $0\le\%6\le1$.
\item $\%7=\rn(1-\%6)\in[0,1]$, since $1-\%6\in[0,1]$ and $0,1\in\Fl$.
\item $\%8=\op{FMA}(\%3,\%7,k)=\rn(\%3\cdot\%7+k)$. Because $\%3=1-k$ exactly, the exact value is
  $k+(1-k)\cdot\%7$, a convex combination of $k$ and $1$, hence in $[k,1]$. By \ref{f:mono} and
  \ref{f:exact}, $\%8\in[k,1]\subseteq[\frac12,1]$.
\item $\%9=\op{FCOPYSIGN}(\%8,x)$. Its obligation, that $x$ is not $\NaN$, holds because $x$ is
  finite. $|\%9|=\%8\in[\frac12,1]$, and $\%9$ carries the sign bit of $x$.
\item $\%10=(\%0\le k)$, decided on non-NaN values, and $\%11=\op{SELECT}(\%10,x,\%9)$.
\item $\%12=\op{FCLAMP}(\%11,-1,1)$. Its obligation, that the bounds are not $\NaN$ and
  $-1\le1$, holds.
\end{itemize}
No port is $\NaN$ and the only partial ops, \op{FCOPYSIGN} and \op{FCLAMP}, have their
obligations discharged, so the kernel does not trap. Two cases remain.
\emph{Case $|x|\le k$.} Then $\%11=x$ with $|x|\le k<1$, the clamp returns $x$ (for $x=-0$,
$\op{FMAX}(-0,-1)=-0$ and $\op{FMIN}(-0,1)=-0$), so $r=x$: it lies in $[-1,1]$ and the sign
clause reads $(x\ge0)\Leftrightarrow(x\ge0)$.
\emph{Case $|x|>k$.} Then $x\neq\pm0$, $\%11=\%9$ with $|\%9|\in[\frac12,1]$, and the clamp
returns $\%9$, so $r\in[-1,1]$. If $x>0$, then $r=\%8>0$ and both sides of the sign clause are
true; if $x<0$, then $r=-\%8<0$ and both are false.
\end{proof}

\begin{remark}[What the proof says about the contract]\label{rem:contract}
The proof uses finiteness of $x$ only to exclude $\NaN$: for $x=\pm\infty$ the same steps give
$\%4=\%5=+\infty$, $\%6=+0$, $\%7=1$, $\%8=1$ and $r=\pm1$, which satisfies the postcondition. The
precondition is therefore stronger than necessary, which is the safe direction. Conversely, the
sign clause does not distinguish $-0$ from $+0$, because IEEE comparison treats them as
equal: the kernel returns $x$ itself on the linear branch, so it preserves a signed zero, but
the contract would also accept a kernel that does not. Stating such intent precisely is the
problem that the project's research track on intent validation addresses.
\end{remark}

\subsection{The first draft fails, and exactly where}

The draft (\file{examples/soft\_knee\_draft.vgs}) differs in two ports:
$\%6=\op{FDIV}(\%4,\%5)$, that is $u=t/(1+t)$, and $\%7=\op{FMIN}(\%6,1)$.

\begin{proposition}[The draft's failure set]\label{prop:draft}
For $\mathit{knee}=\frac12$ and finite $x\in\Fl$, the draft violates its postcondition exactly
when $2^{127}\le|x|\le\Om$. There are $2^{24}=16{,}777{,}216$ such inputs; among them is the
binary32 value of $3\cdot10^{38}$, which is $14791142\cdot2^{104}\approx3.0000000055\times10^{38}$
(\texttt{0x7F61B1E6}). On every violating input the result is $\NaN$, so all three conjuncts of
the postcondition fail.
\end{proposition}
\begin{proof}
With $k=\frac12$, $\%3=\frac12$. If $|x|\le\frac12$, the linear branch returns $x$ and the
postcondition holds as in \cref{thm:softknee}. Let $|x|>\frac12$. If $|x|<2^{24}$, then
$\%1\in(0,2^{24})$ and $\%4=\rn(2\cdot\%1)=2\cdot\%1$ is finite. If $|x|\ge2^{24}$, the spacing of
$\Fl$ at $|x|$ is at least $2$, so $|x|-\frac12$ lies strictly within half a spacing of $|x|$ and
$\%1=\%2=|x|$; then $\%4=\rn(2|x|)$, which by \ref{f:overflow} is $+\infty$ exactly when
$2|x|\ge2^{128}-2^{103}$. The largest element of $\Fl$ below $2^{127}$ is $2^{127}-2^{103}$,
whose double $2^{128}-2^{104}$ does not overflow, and the next element is $2^{127}$, whose double
does. So $\%4=+\infty$ exactly when $|x|\ge2^{127}$.

If $\%4$ is finite, then $\%5=\rn(1+\%4)$ is finite and $\%5\ge\%4$ by \ref{f:mono}, so
$\%6=\rn(\%4/\%5)\in[0,1]$, $\%7=\%6$, $\%8=\rn(\%6/2+\frac12)\in[\frac12,1]$, and the result
$\op{FCOPYSIGN}(\%8,x)$ satisfies the postcondition as in \cref{thm:softknee}. If
$\%4=+\infty$, then $\%5=+\infty$ and $\%6=\infty/\infty=\NaN$; \op{FMIN}, \op{FMA},
\op{FCOPYSIGN} (whose magnitude operand is $\NaN$), \op{SELECT} (which takes $\%9$ because
$|x|>k$) and \op{FCLAMP} all propagate it, so $r=\NaN$ and every comparison in the postcondition
is false.

The elements $x$ of $\Fl$ with $2^{127}\le|x|\le\Om$ are those with biased exponent $254$:
$2^{23}$ of each sign, $2^{24}$ in all. Finally $3\cdot10^{38}/2^{104}\approx14791142.49$ rounds to
$14791142\in[2^{23},2^{24})$, so the binary32 value of $3\cdot10^{38}$ is $14791142\cdot2^{104}$,
which lies in $[2^{127},\Om]$.
\end{proof}

At that input the draft's ports are $\%1=\%2=x$, $\%3=\frac12$, $\%4=\%5=+\infty$ and
$\%6=\dots=\%12=\NaN$, while the repaired kernel computes $\%6=+0$, $\%7=\%8=1$ and $r=1$. This
is the failing slice that the specification's failure report shows the agent.

\subsection{The frame node}

\begin{proposition}[The frame]\label{prop:frame}
If \kw{@soft\_knee} satisfies its contract (\cref{thm:softknee}) and the registry binding for
\kw{@sha256} satisfies \kw{result == spec:sha256(arg0)}, then every obligation of
\kw{@limit\_frame} in \cref{lst:frame} is valid, its cost clause included. (Its \kw{BUDGET}
clause yields no obligation: it is measured.)
\end{proposition}
\begin{proof}
The goal of \op{SIMD\_MAP}, for a fresh index $(i,c)$ with $0\le i<512$ and $0\le c<2$, is the
kernel's precondition at $\mathit{frame}[i,c]$ and $\mathit{knee}$; it follows from the
instance of the node's precondition at $(i,c)$ and from the learnable box. The box's bounds take
the parameter's type \kw{f32}, so they are $\frac12$ and $\kappa$, exactly the bounds that the
kernel's precondition places on $\mathit{knee}$. The first
conjunct of the postcondition, a \kw{forall} over $(i,c)$, becomes a goal at a fresh index; the
fact of \op{SIMD\_MAP} instantiated there is the kernel's postcondition, which states exactly
the range and sign conditions. The second conjunct is the fact of \op{CALL\_HASH}, because
\op{AS\_BYTES}$(\%0)$ denotes $\kw{as\_bytes}(\kw{result.0})$. For the cost clause: the heap is
always empty in the MVP; \op{SIMD\_MAP} over a fixed-size tensor and \op{AS\_BYTES} allocate
nothing, and a callee without the \textsf{ARENA} effect contributes $0$, so
$\kw{arena\_bytes}=0$.
\end{proof}

The proof of the frame is short because the kernel carries the content; the trusted part is the
binding for SHA-256, whose evidence must be \emph{proved} for the frame's postcondition to be a
proof, and the fidelity of \kw{spec:sha256} to FIPS~180-4, which no proof can establish. The
latency budget is measured, never proven. The checked entry point of the exported node
evaluates the precondition, 1,024 finiteness tests, on every external call.

\subsection{Mechanical cross-check}

We checked \cref{thm:softknee,prop:draft} by enumeration. A C program in the paper's directory,
\file{softknee\_check.c}, evaluates both kernels port by port in binary32. IEEE minimum and
maximum are implemented explicitly, \op{FMA} is the C library's \kw{fmaf}, and the program is
compiled with the flags \kw{-O2}, \kw{-ffp-contract=off} and \kw{-fno-fast-math}, so no
operation is fused or reassociated. It also checks the intermediate facts used in the proof: no port is $\NaN$,
$\%2$ is never $-0$, $\%3=1-k$ exactly, the ranges of $\%4$ to $\%8$, and the magnitude and
sign of $\%9$. \Cref{tab:enum} reports a run on an Apple M3 Max under macOS 27.2 (Apple clang
21.0.0, 16 threads, 52\,s of wall time; the program first checks that subnormals are not
flushed).

\begin{table}[t]
\centering
\small
\begin{tabularx}{\linewidth}{@{}l>{\raggedright\arraybackslash}Xrr@{}}
\toprule
& Inputs & Evaluations & Violations \\
\midrule
E1 & every finite $x$ ($4{,}278{,}190{,}080$ values) for 10 values of $k$: $\frac12$, its
  successor, $0.75$, $0.9$, $\kappa$, its predecessor, and 4 seeded random values &
  $42{,}781{,}900{,}800$ & 0 \\
E2 & every $k\in[\frac12,\kappa]$ ($8{,}220{,}837$ values) against 36 special $x$ (signed zeros,
  subnormals, $\pm\Om$, $\pm k$ and its neighbours, \dots) & $295{,}950{,}132$ & 0 \\
E3 & seeded random pairs over the whole precondition & $1{,}073{,}741{,}824$ & 0 \\
E4 & the draft, every finite $x$, $k=\frac12$ & $4{,}278{,}190{,}080$ & $16{,}777{,}216$ \\
\bottomrule
\end{tabularx}
\caption{Exhaustive and random checks of the soft-knee kernels. E1--E3 test the repaired kernel
against its postcondition and the intermediate facts of the proof of \cref{thm:softknee}; E2
also confirms that $1-k$ is exact for every admissible $k$. E4 counts the draft's violations;
the count is exactly $2^{24}$, as \cref{prop:draft} predicts.}
\label{tab:enum}
\end{table}

The draft's failure set grows as the knee rises, because $\%3=1-k$ shrinks and the quotient
overflows sooner: the same run found $33{,}554{,}432$ violating inputs at $k=0.75$ and
$112{,}742{,}912$ at $k=\kappa$. The enumeration checks our reading of the semantics in C, not
\VG's interpreter, which does not exist yet; it confirms the hand proofs on every input it
covers and covers every binary32 $x$ for the knees it lists.
